\equalenv{lem}{lemm}
\equalenv{cor}{coro}
\equalenv{defn}{defi}

\newcommand{\sn}{\mathfrak{S}_n}
\newcommand{\sr}{\mathfrak{S}_r}
\newcommand{\mfs}[1]{\mathfrak{S}_{#1}}
\newcommand{\snn}{\mfs{\smash{[\ol n, n]}}}
\newcommand{\mfsmin}[1]{\mathfrak{S}_-^{#1}}
\newcommand{\bn}{\mathfrak{B}_n}
\newcommand{\br}{\mathfrak{B}_r}
\newcommand{\mfb}[1]{\mathfrak{B}_{#1}}
\newcommand{\mfbmin}[1]{\mathfrak{B}_-^{#1}}
\newcommand{\zsn}{\mathbb{Z}[\sn]}
\newcommand{\qsn}{\mathbb{Q}[\sn]}
\newcommand{\zbn}{\mathbb{Z}[\bn]}
\newcommand{\qbn}{\mathbb{Q}[\bn]}
\newcommand{\cx}{\mathbb{C}[x]}
\newcommand{\zx}{\mathbb{Z}[x]}
\newcommand{\cxn}{\mathbb{C}[x_{1,1},\dotsc,x_{n,n}]}
\newcommand{\zxn}{\mathbb{Z}[x_{1,1},\dotsc,x_{n,n}]}
\newcommand{\zqq}{\mathbb{Z}[q, q^{-1}]}
\newcommand{\csn}{\mathbb{C}[\sn]}
\newcommand{\bitab}{\mathcal B}
\newcommand{\tab}{\mathcal U}
\newcommand{\bctab}{\mathcal U^{\msfBC}}
\newcommand{\atab}{\mathcal U^{\msfA}}
\newcommand{\ol}[1]{\overline{#1}}

\newcommand{\hnq}{H_n(q)}
\newcommand{\hbnq}{H_n^{\msfBC}(q)}
\newcommand{\hbq}[1]{H_{#1}^{\msfBC}(q)}
\newcommand{\hrq}{H_r(q)}
\newcommand{\hlq}{H_\lambda(q)}
\newcommand{\A}{\mathcal{A}}
\newcommand{\anq}{\mathcal{A}(n,q)}
\newcommand{\atwonq}{\mathcal{A}(2n,q)}
\newcommand{\arq}{\mathcal{A}(r;q)}
\newcommand{\arnq}{\mathcal{A}_r}
\newcommand{\annq}{\mathcal{A}_n(n;q)}
\newcommand{\Ann}{\mathcal{A}_{[n],[n]}}
\newcommand{\trsp}{\mathcal{T}}
\newcommand{\btrsp}{\mathcal{T}^{BC}}

\newcommand{\coh}{\mathrm{H}}
\newcommand{\icoh}{\mathrm{IH}}
\newcommand{\HH}{\mathcal{H}}

\newcommand{\redexp}[2]{s_{i_{#1}} \ntnsp \cdots s_{i_{#2}}}
\newcommand{\wtc}[2]{\widetilde{C}_{#1}(#2)}
\newcommand{\btc}[2]{\widetilde{C}_{#1}^{\msfBC}(#2)}

\newcommand{\imm}[1]{\mathrm{Imm}_{#1}}
\newcommand{\bnimm}[1]{\mathrm{Imm}_{#1}^{\bn}}
\newcommand{\simm}[2]{\mathrm{Imm}_{#2}^{\smash{\mfs{#1}}}}

\newcommand{\sumsb}[1]{\sum_{\substack{#1}}}  %%%for multi-line sums
\newcommand{\prodsb}[1]{\prod_{\substack{#1}}}  %%%for multi-line products
\newcommand{\inv}{\textsc{inv}}
\newcommand{\rinv}{\textsc{rinv}}
\newcommand{\defeq}{:=} 
\newcommand{\dfct}{\mathrm{dfct}}
\newcommand{\dnc}{\textsc{dnc}}
\newcommand{\dc}{\textsc{dc}}
\newcommand{\pnc}{\textsc{pnc}}
\newcommand{\sct}{\mathrm{sct}}
\newcommand{\des}{\mathrm{des}}
\newcommand{\exc}{\mathrm{exc}}
\newcommand{\stat}{\mathrm{stat}}
\newcommand{\frobch}{\mathrm{ch}}
\newcommand{\npfrobch}{\mathrm{nch}}
\newcommand{\pfrobch}{\mathrm{pch}}
\newcommand{\spn}{\mathrm{span}}
\newcommand{\sgn}{\mathrm{sgn}}
\newcommand{\triv}{\mathrm{triv}}
\newcommand{\wgt}{\mathrm{wgt}}
\newcommand{\src}{\mathrm{src}}
\newcommand{\snk}{\mathrm{snk}}
\newcommand{\type}{\mathrm{type}}
\newcommand{\ctype}{\mathrm{ctype}}
\newcommand{\inc}{\mathrm{inc}}
\newcommand{\ngr}{\mathrm{ngr}}
\newcommand{\incross}{\textsc{invnc}}
\newcommand{\cdncross}{\textsc{cdnc}}
\newcommand{\cross}{\textsc{c}}
\newcommand{\pavoiding}{$3412$-avoiding, $4231$-avoiding }
\newcommand{\avoidsp}{avoids the patterns $3412$ and $4231${}}
\newcommand{\avoidp}{avoid the patterns $3412$ and $4231${}}
\newcommand{\avoidingp}{avoiding the patterns $3412$ and $4231${}}
\newcommand{\theps}{the patterns $3412$ and $4231${}}
\newcommand{\avoidsignedp}{avoid the signed patterns $1\ol2$, $\ol21$, $\ol2\ol1$, $312$, $3\ol12${}}
\newcommand{\avoidssignedp}{avoids the signed patterns $1\ol2$,  $\ol21$, $\ol2\ol1$, $312$, $3\ol12${}}
\newcommand{\avoidingsignedp}{avoiding the signed patterns $1\ol2$, $\ol21$, $\ol2\ol1$, $312$, $3\ol12${}}
\newcommand{\ul}[1]{\underline{#1}}  
\newcommand{\ssec}[1]{\subsection{#1}{$\negthinspace$}}

\newcommand{\tr}{{\negthickspace \top \negthickspace}}
\newcommand{\nhtnsp}{\kern-0.08333em}
\newcommand{\ntnsp}{\negthinspace}
\newcommand{\ntksp}{\negthickspace}
\newcommand{\nTksp}{\negthickspace\negthickspace}
\newcommand{\nTknsp}{\negthickspace\negthickspace\negthinspace}
\newcommand{\nTtksp}{\negthickspace\negthickspace\negthickspace}
\newcommand{\nTTksp}{\negthickspace\negthickspace\negthickspace
	\negthickspace}
\newcommand{\nTTtksp}{\negthickspace\negthickspace\negthickspace
	\negthickspace\negthickspace}

\newcommand{\oglnc}{{\mathcal O}(\mathrm{GL}_n (\mathbb C))}
\newcommand{\oqglnc}{{\mathcal O}_q(\mathrm{GL}_n (\mathbb C))}
\newcommand{\oslnc}{{\mathcal O}(\mathrm{SL}_n) (\mathbb C)}
\newcommand{\oslthreec}{{\mathcal O}(\mathrm{SL}_3 (\mathbb C))}
\newcommand{\oqslnc}{{\mathcal O}_q(\mathrm{SL}_n (\mathbb C))}
\newcommand{\ospnc}{{\mathcal O}(\mathrm{SP}_{2n} (\mathbb C))}
\newcommand{\oqspnc}{{\mathcal O}_q(\mathrm{SP}_{2n} (\mathbb C))}
\newcommand{\uglnc}{U(\mathfrak{gl}(n,\mathbb{C}))}
\newcommand{\uqglnc}{U_q(\mathfrak{gl}(n,\mathbb{C}))}
\newcommand{\uslnc}{U(\mathfrak{sl}(n,\mathbb{C}))}
\newcommand{\uqslnc}{U_q(\mathfrak{sl}(n,\mathbb{C}))}
\newcommand{\uspnc}{U(\mathfrak{sp}_{2n}(\mathbb{C}))}
\newcommand{\uqspnc}{U_q(\mathfrak{sp}_{2n}(\mathbb{C}))}
\newcommand{\mat}[1]{\mathrm{Mat}_{#1 \times #1}}
\newcommand{\spc}[1]{\mathfrak{sp}_{#1}(\mathbb{C})}
\newcommand{\slc}[1]{\mathfrak{sl}_{#1}(\mathbb{C})}
\newcommand{\fl}[1]{\mathcal{F}_{#1}}
\newcommand{\PiBC}{\Pi^{\msfBC}}


\newcommand{\qdet}{\mathrm{det}_q}
\newcommand{\qper}{\mathrm{per}_q}
\newcommand{\perm}{\mathrm{per}}
\newcommand{\rnk}{\mathrm{rank}}
\newcommand{\sort}{\mathrm{sort}}
\newcommand{\diag}{\mathrm{diag}}
\newcommand{\shape}{\mathrm{shape}}
\newcommand{\hess}{\mathrm{Hess}}
\newcommand{\ahess}{\mathrm{Hess^{\msfA}}}
\newcommand{\bhess}{\mathrm{Hess}^{\msfBB}}
\newcommand{\chess}{\mathrm{Hess}^{\msfC}}
\newcommand{\permmon}[2]{#1_{1,#2_1} \ntnsp\cdots #1_{n,#2_n}}
\newcommand{\bpermmon}[2]
           {#1_{\smash{\ol n,#2_{\ol n}}} \ntnsp\cdots #1_{\smash{\ol 1,#2_{\ol 1}}}
           #1_{1,#2_{1}} \ntnsp\cdots #1_{n,#2_n} }
\newcommand{\sprod}[2]{s_{#1_1} \ntnsp\cdots s_{#1_#2}}
\newcommand{\sintprod}[3]{s_{[#1_1,#2_1]} \ntnsp\cdots s_{[#1_#3, #2_#3]}}
\newcommand{\ssm}{\smallsetminus}
\newcommand{\multiu}{\Cup}
\newcommand{\bs}{\backslash}  
\newcommand{\bfx}{\mathbf x}
\newcommand{\wleq}{\leq_W}
\newcommand{\wless}{<_W}  
\newcommand{\slambda}{\mathfrak{S}_\lambda}
\newcommand{\slambdamin}{\mathfrak{S}_\lambda^{-}}
\newcommand{\phmin}{\phantom{-}}
\newcommand{\circd}[1]{\raisebox{-9pt}{\textcircled{\raisebox{-.9pt}{#1}}}}
\newcommand{\sdot}{\left[ \begin{smallmatrix} 0 & -1 \\ 1 & \phantom{-}0 \end{smallmatrix}\right]}
\newcommand{\sddot}{\left[ \begin{smallmatrix} \phantom{-}0 & 1 \\ -1 & 0 \end{smallmatrix}\right]}
\newcommand{\sreg}{\left[ \begin{smallmatrix} 0 & 1 \\ 1 & 0 \end{smallmatrix}\right]}
\newcommand{\smallj}{\left[ \begin{smallmatrix} 0 & \phantom{\int} & 1 \\ & \smash{\iddots} & \phantom{T} \\ 1 & & 0 \end{smallmatrix}\right]}
\newcommand{\abdiag}[1]{\left[ \begin{smallmatrix} 1 & #1 \\ 0 & 1 \end{smallmatrix}\right]}
\newcommand{\bediag}[1]{\left[ \begin{smallmatrix} 1 & 0 \\ #1 & 1 \end{smallmatrix}\right]}
\newcommand{\upparrow}{\big \uparrow \nTksp \phantom{\uparrow}}
\newcommand{\upparrowa}{\big \uparrow \ntksp \phantom{\uparrow}}
\newcommand{\precdot}{\prec \negthickspace \negthinspace \cdot\ }
\newcommand{\net}[2]{\mathcal F^{\mathsf{#1}}(#2)}
\newcommand{\snet}[2]{\mathcal S^{\mathsf{#1}}(#2)}
\newcommand{\znet}[2]{\mathcal S^{\mathsf{#1}}_{\mathrm Z}(#2)}
\newcommand{\dnet}[2]{\mathcal S^{\mathsf{#1}}_{\mathrm D}(#2)}
\newcommand{\sfA}{\textsf{A}}
\newcommand{\msfA}{\mathsf{A}}
\newcommand{\sfBB}{\textsf{B}}
\newcommand{\msfBB}{\mathsf{B}}
\newcommand{\sfC}{\textsf{C}}
\newcommand{\msfC}{\mathsf{C}}
\newcommand{\sfBC}{\textsf{BC}}
\newcommand{\msfBC}{\mathsf{BC}}
\newcommand{\YBCq}{Y_q^\mathsf{BC}}
\newcommand{\vertex}[1]{#1}
\newcommand{\phsum}{\phantom{\sum_A^Z}}
\newcommand{\phm}{\phantom M}
\newcommand{\phn}{\phantom{ni}}
\newcommand{\bcao}[1]{\mathcal{O}^{\msfBC}({#1})}
\newcommand{\schub}[1]{\Omega_{#1}}
\newcommand{\Aschub}[1]{\Omega_{#1}^{\msfA}}
\newcommand{\Bschub}[1]{\Omega_{#1}^{\msfBB}}
\newcommand{\Cschub}[1]{\Omega_{#1}^{\msfC}}
\newcommand{\tripprox}{\setbox0\hbox{$\approx$}%
\mbox{\makebox[0pt][l]{\raisebox{0.48\ht0}{$\approx$}}$\approx$}}
\newcommand*{\diffeo}{% 
  \mathrel{\vcenter{\offinterlineskip
  \hbox{$\sim$}\vskip-.35ex\hbox{$\sim$}\vskip-.35ex\hbox{$\sim$}}}}

\theoremstyle{remark}

\section{Introduction}\label{s:intro}

Let $W$ be a Coxeter group, $H = H(W)$ its Hecke algebra, and $\trsp(H)$ the
space of Hecke algebra {\em traces},
linear functionals $\theta_q: H \rightarrow \zqq$ satisfying
$\theta_q(DD') = \theta_q(D'D)$ for all $D, D' \in H$.
Included in $\trsp(H)$ are the $H$-characters,
which encode much of the structure of $H$ in a condensed form.
Since traces are linear,
one might hope to solve the
following problem for particular bases
$\mathcal D = \{ D_w \,|\, w \in W \}$ of $H$
and $\Theta = \{ \theta_q^{(i)} \,|\, i = 1,\dotsc,p \}$
of $\trsp(H)$.
\begin{prob}\label{p:evaltrace}
  Find combinatorial formulas for
  each of the trace evaluations
$\{ \theta_q^{(i)} (D_w) \,|\,
\theta_q^{(i)} \ntnsp\in \Theta, D_w \ntnsp\in \mathcal D\}$.
\end{prob}

Unfortunately, trace evaluation is not always easy,
even in type $\msfA$,
when $W$ is the symmetric group $\sn$ with Hecke algebra $H = \hnq$.
(See e.g. \cite[\S 1]{CHSSkanEKL}.)
Type-$\msfA$ solutions were given in
\cite{CSkanTNNChar, KLSBasesQMBIndSgn},
using the induced sign character basis of $\trsp(\hnq)$,
and bases consisting of products of simple elements of
the (modified, signless) Kazhdan--Lusztig basis
$\{ \wtc wq \,|\, w \in \sn \}$ of $\hnq$.
It would be interesting to solve Problem~\ref{p:evaltrace}
for other pairs of type-$\msfA$ bases as well, as these evaluations are related to
facts and conjectures
concerning nonnegativity,
graph coloring, and Hessenberg varieties. (See e.g.
\cite[Lem.\,1.1]{HaimanHecke},
\cite[Conj.\,2.1]{HaimanHecke},  
\cite[Conj.\,4.9]{SWachsChromQ},
\cite[Conj.\,5.5]{StanStemIJT}.)

Partial type-$\msfA$ solutions to Problem~\ref{p:evaltrace}
were given in \cite{CHSSkanEKL, SkanCCS}
for various bases of $\trsp(\hnq)$, and
the subset
\begin{equation}\label{eq:introsmooth}
  \{ \wtc wq \,|\, w \in \sn \text{ \avoidsp} \}
\end{equation}
of the Kazhdan--Lusztig basis of $\hnq$.
By \cite[Thm.\,4.3]{SkanNNDCB},
we have that for $w$ \avoidingp, there exists a planar network $F = F(w)$
which serves as a combinatorial interpretation for $\wtc wq$.
By \cite[Thm.\,7.4]{CHSSkanEKL},
there also exist
a poset $P = P(w)$
and graph $G = G(w)$
such that
evaluations $\theta_q( \wtc wq )$ may be computed combinatorially
by
\begin{enumerate}
\item filling Young diagrams with paths in $F$,
\item filling Young diagrams with elements of $P$,
\item coloring vertices of $G$, 
\item orienting edges of $G$,
\end{enumerate}
while obeying certain rules in each case.
(See also \cite{AthanPSE, SWachsChromQF, StanSymm}.)
While \eqref{eq:introsmooth} is only a subset of the Kazhdan--Lusztig basis
of $\hnq$, it is conjectured~\cite[Conj.\,1.9]{ANigroUpdate},
\cite[Conj.\.3.1]{HaimanHecke} that an even smaller subset
\begin{equation}\label{eq:introcodom}
  \{ \wtc wq \,|\, w \in \sn
  \text{ is codominant, i.e. avoids the pattern } 312 \}
\end{equation}
explains trace evaluations at the entire Kazhdan--Lusztig basis.
It is known~\cite[Thm.\,4.6]{CHSSkanEKL}
that for each element $\wtc wq$ of \eqref{eq:introsmooth}
there exists an element $\wtc vq$ of \eqref{eq:introcodom}
with the property that $P(w) \cong P(v)$
and therefore that $\theta_q(\wtc wq) = \theta_q(\wtc vq)$ for all traces
$\theta_q \in \trsp(\hnq)$.

One could also answer Problem~\ref{p:evaltrace}
from the point of view of symmetric functions.
Let $\Lambda_n$ be the $\mathbb Z$-module of homogeneous, degree-$n$
symmetric functions.
Since the ranks of $\Lambda_n$ and $\trsp(\hnq)$ are
equal,
it is possible to define
a generating function in $\Lambda_n$
for evaluations of traces at
any fixed element $D \in \hnq$.
Following \cite[\S 2]{SkanCCS},
we use
the induced sign character basis $\{\epsilon_q^\lambda\}$ of $\trsp(\hnq)$
and monomial symmetric function basis $\{ m_\lambda \}$ of $\Lambda_n$
to define
\begin{equation*}
  Y_q(D) \defeq \sum_\lambda
  \epsilon_q^\lambda (D) m_\lambda \in \zqq \otimes \Lambda_n.
\end{equation*}
A certain pairing of six natural bases of $\trsp(\hnq)$
and six natural bases of $\Lambda_n$ then guarantees that for each pair
$(\{ \theta_q^\lambda \}, \{ g_\lambda \})$,
we have
\begin{equation*}
  Y_q(D) = \sum_\lambda \theta_q^\lambda(D) g_\lambda.
\end{equation*}
Thus $Y_q(D)$ is in fact a generating function for the evaluation
of all elements of these six trace bases
at $D$. (See e.g. \cite[Prop.\,2.1]{SkanCCS}.)
Conveniently, the combinatorial computations mentioned above also
guarantee
that a certain
{\em chromatic (quasi-)symmetric function}
$X_{G,q}$, defined in terms of the proper colorings of $G$~\cite{SWachsChromQ},
satisfies $X_{G,q} = Y_q(\wtc wq)$~\cite[Thm.\,7.4]{CHSSkanEKL}.
Thus for $w \in \sn$ \avoidingp,
the graph $G$ essentially encodes all trace evaluations of the
form $\theta_q^\lambda(\wtc wq)$ for $\{ \theta_q^\lambda \}$ one of the six natural
bases of $\trsp(\hnq)$.

Some of the above results
from \cite{CHSSkanEKL, SkanNNDCB, SkanCCS}
have type-$\msfBC$ analogs,
i.e. extensions to the hyperoctahedral group $\bn$
and its Hecke algebra $\hbnq$.
In Sections~\ref{s:s2nbn} -- \ref{s:trace},
we present these algebras, their Kazhdan--Lusztig bases,
and their trace spaces.
In Section~\ref{s:planarnet}, we define type-$\msfBC$ analogs
of type-$\msfA$ planar networks, and use these to
graphically represent the subset
\begin{equation}\label{eq:BCintrosmooth}
  \{ \btc w1 \,|\, w \in \bn \text{ \avoidsp} \}
\end{equation}
of the Kazhdan--Lusztig basis of $\zbn$.
In Section~\ref{s:immtnn} we use immanants and total nonnegativity
to interpret trace evaluations at \eqref{eq:BCintrosmooth}
in terms of paths in the type-$\msfBC$ planar networks. 
In Sections~\ref{s:uio} -- \ref{s:incgraph}, we define type-$\msfBC$ analogs
$Q(w)$ and $\Gamma(w)$ of the type-$\msfA$ posets and graphs
associated to planar networks.
We define a type-$\msfBC$ analog of codominant permutations and
show that the posets correspond bijectively to the proper subset
\begin{equation}\label{eq:BCintrocodom}
  \{ \btc w1 \,|\, w \in \bn \text{ $\msfBC$-codominant} \}
\end{equation}
of \eqref{eq:BCintrosmooth}.
We use the above networks, posets, and graphs 
in Section~\ref{s:main}
to state and prove our main results on
the combinatorial computation of type-$\msfBC$ trace evaluations,
and in Section~\ref{s:equiv} to show that
for each element $\btc w1$ of \eqref{eq:BCintrosmooth}
there exists an element $\btc v1$ of \eqref{eq:BCintrocodom}
with the property that $Q(w) \cong Q(v)$
and therefore that $\theta(\btc w1) = \theta(\btc v1)$ for all traces
$\theta \in \trsp(\bn)$.
Formulas in Section~\ref{s:main} lead to natural 
type-$\msfBC$ analogs of
type-$\msfA$ chromatic symmetric functions 
in Section~\ref{s:symm}.
We finish in Section~\ref{s:hess}
with open problems concerning Hessenberg varieties.

\section{The symmetric and hyperoctahedral groups}\label{s:s2nbn}

The hyperoctahedral group $\bn$ is closely related to the symmetric
groups on $n$ and $2n$ letters.  To describe these relationships,
we will use 
subintervals of the set
\begin{equation*}
  [\ol n, n] \defeq \{ -n, \dotsc, n \} \ssm \{ 0 \},
  \end{equation*}
where we define
$\ol a = -a$ for all $a \in [\ol n,n]$.
We call any
subset
$[h,l] \defeq \{ h, \dotsc, l \} \ssm \{0 \}$ of $[\ol n,n]$
an {\em interval}, even if $h < 0 < l$.
Let $\mfs{[h,l]}$ denote
the group of permutations of letters in the interval $[h,l]$.
The group $\bn$ is naturally related both to $\snn$
and $\sn = \mfs{[1,n]}$.
To illustrate these relationships and prepare for our main results,
we will consider the groups' presentations, conjugacy classes, Bruhat
orders, and pattern-avoidance definitions.



\ssec{\texorpdfstring{$\bn$}{B\_n} as a subgroup of \texorpdfstring{$\mfs{\smash{[\ol n, n]}}$}{S\_[n-bar, n]}}\label{ss:bnassubgp}
The $2n$th symmetric group
$\mfs{\smash{[\ol n, n]}}$
is the Coxeter group
(see e.g. \cite{BBCoxeter}) of type $\msfA_{2n-1}$,
with generators
  $s_{\ol{n-1}}, \dotsc, s_{\ol 1}, s_0, s_1, \dotsc, s_{n-1}$
and relations
\begin{equation*}
  \begin{alignedat}2
    s_i^2 &= e &\quad &\text{for $i = \ol{n-1}, \dotsc, n-1$,}\\
    s_is_j &= s_js_i &\quad &\text{for $|i-j| \geq 2$,}\\
    s_is_js_i &= s_js_is_j &\quad &\text{for $|i-j| = 1$.}
  \end{alignedat}
\end{equation*}

If an expression $s_{i_1} \cdots s_{i_\ell}$ for $w \in \snn$
is as short as possible,
then call it {\em reduced} and call $\ell = \ell(w)$ the {\em length} of $w$.
Define a (left) action of $\snn$ on rearrangements
$w_{\ol n} \cdots w_{\ol 1} w_1 \cdots w_n$ of the word
$\ol n \cdots \ol 1 1 \cdots n$ by
\begin{equation}\label{eq:laction}
  \begin{cases}
    s_i \text{ swaps letters in positions $i, i+1$}
      &\text{for $i = 1, \dotsc, n-1$},\\
    s_{\ol i} \text{ swaps letters in positions $\ol i, \ol{i+1}$}
      &\text{for $i = 1, \dotsc, n-1$},\\
    s_0 \text{ swaps letters in positions $\ol 1, 1$,}
  \end{cases}
\end{equation}
and
define the {\em one-line notation} of $w = s_{i_1} \cdots s_{i_r} \in \snn$ to be
\begin{equation}\label{eq:s2nrearrange}
  w_{\ol n} \cdots w_{\ol 1} w_1 \cdots w_n 
= s_{i_1}(s_{i_2}( \cdots (s_{i_r}( \ol n \cdots \ol 1 1 \cdots n)) \cdots )).
\end{equation}
For example, when $n=4$, the element $s_{\ol1}s_0s_1$
has one-line notation
\begin{equation*}
    s_{\ol1}(s_0( s_1(\ol4 \ol3 \ol2 \ol1 1234)))
    = s_{\ol1}(s_0(\ol4 \ol3 \ol2 \ol1 2 1 3 4))
    = s_{\ol1}(\ol4 \ol3 \ol2 2 \ol1 1 3 4)
    = \ol4\ol3 2 \ol2 \ol1 1 3 4.
  \end{equation*}
(By our definition, the right action of $s_i$ swaps the letters $i, i+1$,
wherever they are.)
It follows that $w_i^{-1}$ is the index $j$ satisfying $w_j = i$.
It is known that $\ell(w)$ equals the number of {\em inversions} in $w$:
\begin{equation*}
  \inv(w_{\ol n} \cdots w_{\ol1} w_1 \cdots w_n) \defeq
  \{ (j,i) \,|\, j \mkern -2mu > \mkern -2mu i \text{ and $j$ appears before $i$ in $w_{\ol n} \cdots w_{\ol1} w_1 \cdots w_n$} \}.
  \end{equation*}
Thus we have
  $\ell(\ol4\ol32\ol2\ol1134) =
  \inv(\ol4\ol32\ol2\ol1134) =
  |\{(2,\ol2), (2,\ol1), (2,1)\}| = 3$.

For $[a, b] \subseteq [h, l]$
let $\smash{s_{[a,b]}^{[h,l]}}$ be the permutation
in $\mfs{[h,l]}$ whose one-line notation has the form
\begin{equation*}
  h \cdots (a-1) \cdot b (b-1) \cdots (a+1) a \cdot (b+1) \cdots l.
\end{equation*}
When the interval $[h,l]$ is clear from context,
we will simply write $s_{[a,b]}$.
Call such an element a (type-$\msfA$) {\em reversal}.
Observe that the standard generators of $\snn$
are all reversals:
  $s_0 = s_{[-1,1]}$,
  and for $i \geq 1$ we have
  $s_i = s_{[i,i+1]}$,
  $s_{\smash{\ol i}} = s_{\smash{[\ol{i+1},\ol i]}}$. 
Also observe that each trivial reversal $s_{[a,a]}$
is equal to the identity element $e$, and that two reversals $s_{[a,b]}$,
$s_{[c,d]}$ commute
if their intervals $[a,b]$, $[c,d]$ do not intersect.
Let $\bn$ be the Coxeter group of type $\mathsf C_n = \mathsf B_n$,
i.e. the hyperoctahedral group.
We may view $\bn$ as the subgroup of $\snn$
generated by elements
\begin{equation*}
  t\ (\text{also written $s'_0$}) \defeq s_0, \qquad
  s'_i \defeq s_is_{\smash{\ol i}}, \text{ for } i = 1,\dotsc,n-1,
\end{equation*}
which satisfy the relations
\begin{equation*}
  \begin{alignedat}2
    \smash{{s'_i}^2} &= e &\quad &\text{for $i = 0, \dotsc, n-1$,}\\
    ts'_1ts'_1 &= s'_1ts'_1t, &\quad & \\
    s'_is'_j &= s'_js'_i &\quad &\text{for $i,j \geq 0$ and $|i-j| \geq 2$,}\\
    s'_is'_js'_i &= s'_js'_is'_j &\quad &\text{for $i,j \geq 1$ and $|i-j| = 1$.}
  \end{alignedat}
\end{equation*}
The one-line notation for elements of $\bn$ is inherited from that
of $\snn$.
For example, when $n=4$, the element $ts'_1s'_2 \in \mfb4 $
has one-line notation
\begin{equation*}
    t(s'_1(s'_2(\ol4 \ol3 \ol2 \ol1 1234)))
    = ts'_1(\ol4 \ol2 \ol3 \ol1 1 3 2 4))
    = t(\ol4 \ol2 \ol1 \ol3 3 1 2 4)
    = \ol4 \ol2 \ol1 3 \ol3 1 2 4.
  \end{equation*}
If an expression $s'_{i_1} \cdots s'_{i_\ell}$ for $w \in \bn$
is as short as possible,
then call it {\em reduced} and call $\ell = \ell(w)$ the {\em length} of $w$.
Let $\ell_t(w)$ be the number of occurrences of $t$ in any (equivalently, every)
reduced expression for $w$.
Analogously, let $\ell_s(w) = \ell(w) - \ell_t(w)$ be the number
of occurrences of $s'_1,\dotsc,s'_{n-1}$.
It is easy to see that one-line notations of
elements of $\bn$ are precisely the set of
permutations $w_{\ol n} \cdots w_{\ol 1} w_1 \cdots w_n \in \snn$
which satisfy $w_{\ol i} = \ol{w_i}$,
i.e. each is completely determined by the $n$-letter subword
$w_1 \cdots w_n$.
Call these words the {\em long} and
{\em short}
one-line notations of $w$, respectively.
We can read $\ell(w)$, $\ell_t(w)$, $\ell_s(w)$ from the
short one-line notation of $w$ by
\begin{equation*}%\label{eq:ellandellt}
  \ell(w) = \inv(w_1 \cdots w_n) + \sumsb{i > 0\\ w_i < 0} |w_i|,
  \qquad
  \ell_t(w) = \# \{ i > 0 \,|\, w_i < 0 \},
  \end{equation*}
and $\ell_s(w) = \ell(w) - \ell_t(w)$.
Thus we have
  $\ell(\ol3124) = |3| = 3$, $\ell_t(\ol3124) = 1$, and $\ell_s(\ol3124) = 2$.

Define {\em type-$\msfBC$ reversals} to be those elements of $\bn$
having the forms
\begin{equation}\label{eq:creversal}
  \begin{aligned}
    s'_{\smash{[\ol a, a]}} &\defeq s_{\smash{[\ol a, a]}},
        \mbox{ for } 1 \leq a \leq n,\\
    s'_{\smash{[a,b]}} &\defeq s_{\smash{[\ol b, \ol a]}} s_{[a,b]},
    \mbox{ for } 1 \leq a \leq b \leq n,
  \end{aligned}
\end{equation}
where the elements $s_{[a,b]}$ are reversals in $\snn$.

\ssec{Conjugacy classes, partitions, tableaux, bipartitions, bitableaux}\label{ss:conjclasses}
Conjugacy classes of $\sn$ correspond to {\em (integer) partitions} of $n$,
weakly decreasing positive integer sequences
$\lambda = (\lambda_1, \dotsc, \lambda_\ell)$
satisfying $\lambda_1 + \cdots + \lambda_\ell = n$.
The $\ell = \ell(\lambda)$ components of a partition $\lambda$
are called its {\em parts}
and we let the expressions
$|\lambda| = n$ and $\lambda \vdash n$ denote that $\lambda$ is
a partition of $n$.
Sometimes we use the notation $k^{a_k}$ to denote a
sequence of $a_k$ copies of the letter $k$.
Given $\lambda \vdash n$, we define the {\em transpose}
partition $\lambda^\tr = (\lambda_1^\tr, \dotsc,\lambda_{\lambda_1}^{\tr})$ by
$\lambda_i^\tr = \# \{ j \,|\, \lambda_j \geq i \}$. Thus $n^\tr = 1^n$.
We call $\lambda$ {\em self-transpose} if $\lambda^\tr = \lambda$
and we define the empty sequence $\emptyset$ to be the unique partition of the
integer $0$.
Generalizing integer partitions of $n$ are
{\em compositions} $\alpha = (\alpha_1,\dotsc, \alpha_r)$ of $n$,
which are simply positive integer sequences summing to $n$.
We let the notation $\alpha \vDash n$ denote that $\alpha$ is
a composition of $n$. (See e.g. \cite[\S 1.2]{StanEC1}.)

The conjugacy class of $\sn$ corresponding to $\lambda \vdash n$
is the set of all permutations having cycle type $\lambda$.  We write
$\ctype(w) = \lambda$. Letting
$a_k$ be the multiplicity of $k$ in $\lambda$, for $k =1,\dotsc,n$,
we may express the cardinality of the $\lambda$-conjugacy class of $\sn$ as
$n!/z_\lambda$, where
\begin{equation}\label{eq:zlambda}
 z_\lambda = {1^{a_1} \cdots n^{a_n} a_1! \cdots a_n!}.
\end{equation}

Conjugacy classes of $\bn$ correspond to
{\em integer bipartitions of $n$}, pairs $(\lambda,\mu)$
of integer partitions with 
$|\lambda| + |\mu| = n$.
We let
$(\lambda,\mu) \vdash n$ denote that $(\lambda,\mu)$ is
a bipartition of $n$. 
To explicitly describe the conjugacy classes of $\bn$,
we define the homomorphism
\begin{equation}\label{eq:varphi}
  \begin{aligned}
    \varphi: \bn &\rightarrow \sn \\
    s'_i &\mapsto s_i, \quad i=1,\dotsc,n-1,\\
    t &\mapsto e,
  \end{aligned}
  \end{equation}
which replaces letters in the short one-line notation of
$v \in \bn$ by their absolute values.
For each element $v \in \bn$
and each cycle $C = (c_1, \dotsc, c_k = c_0)$
of $\varphi(v) \in \sn$,
define the {\em signed cycle}
$\hat C = (\hat c_1, \dotsc, \hat c_k)$
of $v$ by
\begin{equation*}
  \hat c_i = \begin{cases}
  c_i &\text{if $v(c_{i-1}) = c_i$},\\
  \ol c_i &\text{if $v(c_{i-1}) = \ol{c_i}$}.
  \end{cases}
\end{equation*}
Call $\hat C$ {\em positive} if it has an even number of negative letters,
and {\em negative} otherwise.
The conjugacy class of $\bn$ corresponding to $(\lambda,\mu) \vdash n$
is precisely the set of elements  
whose {\em signed cycle type}, the bipartition of
positive cycle cardinalities and negative cycle cardinalities,
is equal to $(\lambda,\mu)$.
We write $\sct(w) = (\lambda,\mu)$.
We may express the cardinality of the $(\lambda,\mu)$
conjugacy class of $\bn$ as $2^nn!/(z_\lambda z_\mu 2^{\ell(\lambda) + \ell(\mu)})$.
  

To each integer partition
$\lambda = (\lambda_1, \dotsc, \lambda_\ell) \vdash n$
we associate a {\em Young diagram of shape $\lambda$},
an arrangement of $n$ boxes into $\ell$ left-justified rows
with $\lambda_i$ boxes in row $i$.  By the French convention, row $1$ appears
on the bottom.
A Young diagram filled with
elements of a set $S$ is called a {\em tableau} or more specifically an
{\em $S$-tableau}.  If $S \subseteq \mathbb Z$ we also call it a
{\em Young tableau}.
Repeated elements are permitted.
Given a bipartition $(\lambda,\mu) \vdash n$, we define
a Young {\em bidiagram} of shape $(\lambda,\mu)$
to be an ordered pair of Young diagrams of shapes $\lambda$ and $\mu$.
We define {\em bitableaux} similarly.


\ssec{The Bruhat order}

For any Coxeter group $W$,
the {\em Bruhat order on $W$} is the poset
defined by declaring $v \leq_W w$ 
if
some (equivalently, every)
reduced expression for $w$ contains a reduced expression for $v$.
Ehresmann~\cite{Ehresmann}
showed that the Bruhat order on $\snn$ is isomorphic
to the (dual of the)
componentwise order on tableaux $\{ A(w) \,|\, w \in \snn \}$
of shape $(2n,2n-1,\dotsc,1)$ defined by placing the
increasing rearrangement of $w_i \cdots w_n$ in row $i$,
for $i = \ol n, \dotsc, n$.
